\section{在线派单：覆盖率、平均与最差体验}
\subsection{事件、阈值与结束行为}
Worker和Request均有唯一ID、坐标和到达时刻；同刻先处理请求、再处理司机，各自按ID稳定排序。BB在同刻全部到达后才执行该时刻的批次；结束时补最后非重复批次。派单不早于双方可用，接驾可以在仿真horizon之后完成。对请求$r$和司机$w$，
\begin{equation}
D_{wr}=(t_{assign}-t_r)+\frac{\lVert x_w-x_r\rVert_2}{v}.
\end{equation}
第一项是乘客已等时间，第二项是接驾行驶时间；分别保存，不用单一总数隐藏机制。

若只考虑当前接驾距离，可设计“在已到达且可行的司机中取最近者”的在线贪心；对一个已知批次则用最小总距离匹配优化整体和，若关心最差体验就改求瓶颈。最近的局部选择不保证未来全局最优，这也是继续比较TBG和批次策略的原因。

TBG的$\gamma$是时间阈值。请求到达时，在已到司机中取行驶时间不超过$\gamma$的最近者，否则缓冲。司机到达时先服务FIFO首个等待$\ge\gamma$的请求，不再限制距离；没有超时请求才在总延迟$<\gamma$的候选中取最小者，否则司机缓冲。该行为按本轮提示词和Notebook实现；课件伪码的严格/非严格边界与Notebook有差异，不能悄悄混合。结束兜底只匹配真实剩余双方，供给不足明确保留未服务对象。

BB每隔$\theta$对缓冲双方匹配。Notebook的线性分配最小化总延迟，课件另提最小最大延迟，故分别标为BBsum和BBmax。正式版本回收未匹配对象；原starter清空双方的discard策略作为单独对照。半径或容量等约束可通过共同可行边与顶点容量表达；没有可行边就不匹配，不能强行连边凑满服务数。

\subsection{用什么离线参考计算实例比值}
完全知道未来时，单次匹配每条边的最早派单为$\max(t_w,t_r)$，成本$C_{wr}=\max(t_w-t_r,0)+d_{wr}/v$。先固定最大可行基数$k$，再分别求$\min\sum C$和$\min\max C$。前者使用线性分配\cite{scipy}，后者对边成本阈值做最大基数可行性检查。

例如$C=\left[\begin{smallmatrix}0&6\\6&10\end{smallmatrix}\right]$，sum最优取对角线，总10、最大10；瓶颈最优取非对角线，总12、最大6。调用KM不能自动证明最差延迟最优。小图枚举、独立线性规划与另写二部可行性程序核对了两种参考。

在线目标值除以同实例、同基数、同目标的离线值，才得到本报告的实例比值。零分母、无匹配和基数不同均保留不可计算状态，不能填0或把不完整服务与全服务最优直接相除。有限输入上的实测比值也不构成全输入空间的理论竞争保证。

\subsection{丢单为何会制造较好的条件延迟}
\begin{table}[htbp]\centering\small
\caption{五个开发实例：覆盖和延迟必须一起解释。}\label{tab:onlinedevelopment}
\begin{tabularx}{\linewidth}{Yrrrrr}\toprule
方法(参数) & 匹配率 & 均值/min & 最大/min & P95/min & max实例比\\\midrule
TBG(2) & 1.00 & 7.399 & 13.342 & 11.876 & 1.649\\
TBG(5) & 1.00 & 5.119 & 12.802 & 12.119 & 1.568\\
BBsum(1) & 1.00 & 5.748 & 15.426 & 12.157 & 1.872\\
BBmax(1) & 1.00 & 5.965 & 14.598 & 12.023 & 1.794\\
discard(0.5) & 0.25 & 4.379 & 6.516 & 6.316 & --\\
离线sum & 1.00 & 4.103 & 11.487 & 9.412 & 1.406\\
离线max & 1.00 & 5.747 & 8.370 & 8.145 & 1.000\\
\bottomrule\end{tabularx}
\tabnote{每行对5个输入seed先逐实例算指标，再平均；P95不是把全部seed混成单条观测。部分覆盖的实例比无定义，以--表示。}
\end{table}

教师20×20场景的五个开发seed中，$\theta=0.5$的discard平均覆盖25\%，回收版本覆盖全部请求；司机先到而请求后到的场景中，discard甚至完全没有服务。因此图\ref{fig:onlinecdf}以全部请求为分母，未匹配质量留在较低平台。另一半图用每个实例的平均与最大延迟成对展示：改善均值不自动改善尾部。

\datafigure{online_delay_and_coverage}{教师物理场景，5个成对输入seed、每seed20请求。左侧每点是一个实例，十字为跨seed均值；右侧分母100请求，黑线为discard，覆盖25\%。}{fig:onlinecdf}

\subsection{阈值、到达顺序与局部瓶颈的边界}
原始控制变量包含供需比、热点、司机/请求先到、规模与批长；同一因素变化以相同输入和成对seed比较。教师场景中TBG小阈值0.5/1/2均有平均17\%服务推迟到结束兜底，阈值5下降到7\%，但阈值继续增大并不保证更低最大延迟。由这一观察触发3种到达顺序下的局部阈值实验（图\ref{fig:gamma}），而非预先堆叠增强。

\datafigure{online_threshold_diagnosis}{阈值的针对性开发实验。每点为5个相同输入seed的均值，误差条为seed标准差；右侧是结束兜底比例。线只连接实测阈值。}{fig:gamma}
\datafigure{online_spatial_mechanism}{同一教师实例的空间匹配与等待/行驶分解。案例按TBG5最大延迟中位seed选择；连线是接驾边，不代表真实道路。每个面板均使用同一20对输入。}{fig:spatial}

BBmax在一个批次内最小化最大成本，但会改变留给未来的司机。开发teacher seed13中，$\theta=1$的BBmax全天最大延迟反而比BBsum高2.162 min，不能把局部数学最优扩大为整个在线过程保证。供需不平衡和到达顺序也改变阈值的作用（图\ref{fig:scenes}）。没有对这些负例改写指标。

\datafigure{online_scenario_matrix}{6种真实运行场景、5seed均值，颜色表示同基数的最大延迟实例比值，越小越好。每格是实测设置，没有插值。}{fig:scenes}
\begin{table}[htbp]\centering\small
\caption{五个独立合成确认实例：不按结果重新选择阈值。}\label{tab:onlineconfirmation}
\begin{tabularx}{\linewidth}{Yrrrrr}\toprule
方法(参数) & 匹配率 & 均值/min & 最大/min & P95/min & max实例比\\\midrule
TBG(2) & 1.00 & 7.639 & 13.581 & 13.007 & 1.936\\
TBG(5) & 1.00 & 5.411 & 15.927 & 14.070 & 2.272\\
BBsum(1) & 1.00 & 5.255 & 12.208 & 10.111 & 1.727\\
BBmax(1) & 1.00 & 5.369 & 12.208 & 10.291 & 1.727\\
discard(0.5) & 0.37 & 5.348 & 9.239 & 8.808 & --\\
离线sum & 1.00 & 4.033 & 10.645 & 7.863 & 1.497\\
离线max & 1.00 & 5.136 & 7.050 & 6.914 & 1.000\\
\bottomrule\end{tabularx}
\tabnote{每行对5个输入seed先逐实例算指标，再平均；P95不是把全部seed混成单条观测。部分覆盖的实例比无定义，以--表示。}
\end{table}

独立合成确认使用seed110--114，没有据其结果再选阈值。TBG5的平均延迟5.411 min低于TBG2的7.639 min，但平均最大延迟15.927 min高于13.581 min：开发样本中阈值5对尾部的轻微优势没有在确认样本保持。因此不把阈值5推荐为普适改进；平均与最差体验仍需分别判断。结论只适用于所列到达和空间模型。

\subsection{真实签到对空间需求的扩展}
从inner训练的4个小时池重采样20个代理起点，形成各自60 min流；司机位置也只取训练签到，司机到达时间为模拟值。空间采用原条件工作坐标、速度1 km/min，与10×10 km均匀教学场景分开标识。图\ref{fig:proxyhours}比较时段和方法，不把重采样批次称为真实运营小时。
\datafigure{online_proxy_hours}{8/13/18/23点来源的历史代理起点，各5个采样seed。左、右分别为已匹配平均延迟和最大延迟，20司机/20请求；教学局部km与模拟分钟。}{fig:proxyhours}
